\section{Introduction}
\label{sec:intro}

A common way to show a large image or document collection is to embed the items in
the plane (for example with t-SNE~\citep{vandermaaten2008tsne} or
UMAP~\citep{mcinnes2018umap}) and then snap every item to its own cell of a regular grid,
so that thumbnails tile the screen without overlap. Snapping is an assignment problem:
each point must receive a distinct cell, and we would like points to move as little as
possible. Measured as total squared displacement, the best snapping is a linear
assignment problem (LAP), and also a discrete optimal-transport problem with quadratic
cost~\citep{peyre2019computational}.

For a few thousand points, practitioners solve this LAP exactly with a dense
Jonker--Volgenant or Hungarian solver~\citep{markovtsev2017jvtsne,mcdonald2016cloudtogrid,
markovtsev_lapjv,kazmar_lap}. The dense cost matrix needs $8N^2$ bytes, which is
80\,GB at $N=100{,}000$, so these solvers stop being usable somewhere between
10{,}000 and 30{,}000 points. Beyond that, people turn to recursive-splitting heuristics
such as RasterFairy~\citep{klingemann2015rasterfairy} or to fast approximate transport
such as BSP-OT~\citep{genest2025bspot}.

This paper develops and evaluates \ShuffleSnap, a simple heuristic for this regime,
taking its name from Shellsort~\citep{shell1959highspeed}. Shellsort sorts interleaved
subsequences at a decreasing sequence of gaps. \ShuffleSnap\ does the same in 2D: at stride $s$
it splits the grid into $s^2$ interleaved subgrids and solves small \emph{exact}
assignment problems inside $6\times6$ windows of each subgrid, halving $s$ until it
reaches 1 (Figure~\ref{fig:method}). All work is at full resolution; nothing is
aggregated or downsampled.

We treat the idea as a hypothesis to test, not a result to promote. Our aim throughout
was to find where it fails. The main contributions are:

\begin{itemize}
\item \textbf{A careful implementation} of the specified baseline, with documented
  boundary, rectangular-grid, surplus-cell and mask rules. It is a parallel C++ kernel
  with a Python API whose results do not depend on the number of threads.
\item \textbf{A certified exact solver that scales to $10^5$ points.} It solves min-cost
  flow on sparse candidate edges, then prices every edge of the \emph{dense} problem
  through a Lagrangian bound, adding edges until the bound meets the primal cost. It
  gives true optimality gaps, not progress relative to our own best answer, for every
  instance up to 200{,}000 points in this study. In our experiments it is also a strong
  practical baseline.
\item \textbf{Measured improvements and non-improvements.} On a held-out development set,
  a finer stride sequence (ratio $\sqrt2$ instead of 2) and larger windows reduce the
  excess cost severalfold for a small time increase. Several intuitive ideas make things
  \emph{worse}: sorted or bisection starts, polishing at coarser strides, and a
  Jonker--Volgenant window solver.
\item \textbf{A diagnosis of the remaining error.} The gap to the optimum is concentrated
  in a few long closed loops of points, each displaced by about one cell. No window
  smaller than a loop can remove it. Masks with concave bays and grids with surplus
  cells expose the same mechanism at a larger scale.
\item \textbf{A benchmark against exact LAP, RasterFairy and BSP-OT.} We first verified
  each comparator's capabilities: we found and worked around coordinate truncation in the
  PyPI release of RasterFairy, and used its much stronger upstream version. The
  benchmark covers varied and adversarial distributions, masks, prime point counts, two
  real UMAP embeddings with 70{,}000 points each, multiple seeds, matched-budget curves,
  and scaling to $10^6$ points. We report cost, time, peak memory, and neighbourhood
  preservation, and we keep displacement and neighbourhood measures separate.
\end{itemize}

The implementation (\texttt{shufflesnap}) and this paper's experiment code, raw results
and figure scripts are released as two separate repositories.
